\documentclass[../../main.tex]{subfiles}
\begin{document}
% SOURCE p=218; printed=203; object=ch06-marcinkiewicz; number=6.28; region=upper
\begin{theorem}[Marcinkiewicz 插值定理]\label{theorem:ch06-marcinkiewicz}
设 $(X,\mathscr M,\mu)$ 与 $(Y,\mathscr N,\nu)$ 为测度空间；$p_0,p_1,q_0,q_1\in[1,\infty]$，满足 $p_0\leqslant q_0$、$p_1\leqslant q_1$、$q_0\ne q_1$；并设
\begin{align*}
\frac1p=\frac{1-t}{p_0}+\frac t{p_1},\qquad
\frac1q=\frac{1-t}{q_0}+\frac t{q_1},\qquad 0<t<1.
\end{align*}
若 $T$ 为从 $L^{p_0}(\mu)+L^{p_1}(\mu)$ 到 $Y$ 上可测函数空间的次线性映射，且同时为弱 $(p_0,q_0)$ 型与弱 $(p_1,q_1)$ 型，则 $T$ 为强 $(p,q)$ 型。更准确地，若 $[Tf]_{q_j}\leqslant C_j\|f\|_{p_j}$，$j=0,1$，则
\begin{align*}
\|Tf\|_q\leqslant B_p\|f\|_p,
\end{align*}
其中 $B_p$ 除了 $p$ 之外，只依赖 $p_j,q_j,C_j$；当 $p\to p_j$ 时，若 $p_j<\infty$，则 $B_p|p-p_j|$ 保持有界，若 $p_j=\infty$，则 $B_p$ 保持有界。
\end{theorem}
% SOURCE p=218,219,220; printed=203-205; object=ch06-marcinkiewicz-proof; region=middle-bottom
\begin{proof}
$p_0=p_1$ 的情形很容易，原书留给读者，见\cref{exercise:ch06-marcinkiewicz-equal-input}。故不妨设 $p_0<p_1$，暂且再设 $q_0,q_1<\infty$，于是 $p_0<p_1<\infty$。给定 $f\in L^p(\mu)$、$A>0$，按\cref{proposition:ch06-distribution-truncation}定义 $g_A,h_A$。由层积分公式及其分布函数，
\begin{align}
\int|g_A|^{p_0}\,\mathrm d\mu
&=p_0\int_0^\infty\beta^{p_0-1}\lambda_{g_A}(\beta)\,\mathrm d\beta
=p_0\int_0^\infty\beta^{p_0-1}\lambda_f(\beta+A)\,\mathrm d\beta\notag\\
&=p_0\int_A^\infty(\beta-A)^{p_0-1}\lambda_f(\beta)\,\mathrm d\beta
\leqslant p_0\int_A^\infty\beta^{p_0-1}\lambda_f(\beta)\,\mathrm d\beta,\notag\\
\int|h_A|^{p_1}\,\mathrm d\mu
&=p_1\int_0^\infty\beta^{p_1-1}\lambda_{h_A}(\beta)\,\mathrm d\beta
=p_1\int_0^A\beta^{p_1-1}\lambda_f(\beta)\,\mathrm d\beta.
\label{align:ch06-marcinkiewicz-input-truncations}
\end{align}
同样，
\begin{align}
\int|Tf|^q\,\mathrm d\nu
=q\int_0^\infty\alpha^{q-1}\lambda_{Tf}(\alpha)\,\mathrm d\alpha
=2^qq\int_0^\infty\alpha^{q-1}\lambda_{Tf}(2\alpha)\,\mathrm d\alpha.
\label{align:ch06-marcinkiewicz-output-layers}
\end{align}
因为 $T$ 次线性，由\subcref{proposition:ch06-distribution-properties}{item:ch06-distribution-sum}，
\begin{align}
\lambda_{Tf}(2\alpha)\leqslant\lambda_{Tg_A}(\alpha)+\lambda_{Th_A}(\alpha).
\label{align:ch06-marcinkiewicz-output-split}
\end{align}
此式对所有 $\alpha,A>0$ 成立，故可以让 $A$ 依赖 $\alpha$。由定义 $p,q$ 的插值关系，
\begin{align}
\frac{p_0(q_0-q)}{q_0(p_0-p)}
&=\frac{p^{-1}(q^{-1}-q_0^{-1})}{q^{-1}(p^{-1}-p_0^{-1})}
=\frac{p^{-1}(q^{-1}-q_1^{-1})}{q^{-1}(p^{-1}-p_1^{-1})}\notag\\
&=\frac{p_1(q_1-q)}{q_1(p_1-p)}.
\label{align:ch06-marcinkiewicz-sigma}
\end{align}
记共同值为 $\sigma$，取 $A=\alpha^\sigma$。由上述公式及端点弱型估计，
\begin{align}
\|Tf\|_q^q
&\leqslant2^qq\int_0^\infty\alpha^{q-1}
\left[\left(\frac{C_0\|g_A\|_{p_0}}\alpha\right)^{q_0}
+\left(\frac{C_1\|h_A\|_{p_1}}\alpha\right)^{q_1}\right]\mathrm d\alpha\notag\\
&\leqslant2^qqC_0^{q_0}p_0^{q_0/p_0}\int_0^\infty\alpha^{q-q_0-1}
\left[\int_{\alpha^\sigma}^\infty\beta^{p_0-1}\lambda_f(\beta)\,\mathrm d\beta\right]^{q_0/p_0}\mathrm d\alpha\notag\\
&\quad+2^qqC_1^{q_1}p_1^{q_1/p_1}\int_0^\infty\alpha^{q-q_1-1}
\left[\int_0^{\alpha^\sigma}\beta^{p_1-1}\lambda_f(\beta)\,\mathrm d\beta\right]^{q_1/p_1}\mathrm d\alpha\notag\\
&=\sum_{j=0}^1 2^qqC_j^{q_j}p_j^{q_j/p_j}
\int_0^\infty\left[\int_0^\infty\phi_j(\alpha,\beta)\,\mathrm d\beta\right]^{q_j/p_j}\mathrm d\alpha.
\label{align:ch06-marcinkiewicz-two-integrals}
\end{align}
这里 $\chi_0,\chi_1$ 分别是 $\{(\alpha,\beta):\beta>\alpha^\sigma\}$、$\{(\alpha,\beta):\beta<\alpha^\sigma\}$ 的示性函数，并且
\begin{align*}
\phi_j(\alpha,\beta)=\chi_j(\alpha,\beta)
\alpha^{(q-q_j-1)p_j/q_j}\beta^{p_j-1}\lambda_f(\beta).
\end{align*}
原书上式最后一个求和行的系数印作 $C_j^{p_j}$；与前两行及后续常数式一致的系数应为这里明确写出的 $C_j^{q_j}$，其核对说明见后注。

由于 $q_j/p_j\geqslant1$，可应用\cref{theorem:ch06-integral-minkowski}，得到
\begin{align}
\int_0^\infty\left[\int_0^\infty\phi_j(\alpha,\beta)\,\mathrm d\beta\right]^{q_j/p_j}\mathrm d\alpha
\leqslant\left\{\int_0^\infty\left[\int_0^\infty\phi_j(\alpha,\beta)^{q_j/p_j}\,\mathrm d\alpha\right]^{p_j/q_j}\mathrm d\beta\right\}^{q_j/p_j}.
\label{align:ch06-marcinkiewicz-minkowski}
\end{align}
令 $\tau=1/\sigma$。若 $q_1>q_0$，则 $q-q_0>0$、$\sigma>0$，且 $\beta>\alpha^\sigma$ 等价于 $\alpha<\beta^\tau$。于是
\begin{align*}
\int_0^\infty\left[\int_0^\infty\phi_0^{q_0/p_0}\,\mathrm d\alpha\right]^{p_0/q_0}\mathrm d\beta
&=\int_0^\infty\left[\int_0^{\beta^\tau}\alpha^{q-q_0-1}\,\mathrm d\alpha\right]^{p_0/q_0}
\beta^{p_0-1}\lambda_f(\beta)\,\mathrm d\beta\\
&=(q-q_0)^{-p_0/q_0}\int_0^\infty
\beta^{p_0-1+\tau p_0(q-q_0)/q_0}\lambda_f(\beta)\,\mathrm d\beta\\
&=|q-q_0|^{-p_0/q_0}\int_0^\infty\beta^{p-1}\lambda_f(\beta)\,\mathrm d\beta\\
&=|q-q_0|^{-p_0/q_0}p^{-1}\|f\|_p^p.
\end{align*}
第三行用\cref{align:ch06-marcinkiewicz-sigma}化简指数。另一方面，若 $q_1<q_0$，则 $q-q_0<0$、$\sigma<0$，原区域等价于 $\alpha>\beta^\tau$，所以
\begin{align*}
\int_0^\infty\left[\int_0^\infty\phi_0^{q_0/p_0}\,\mathrm d\alpha\right]^{p_0/q_0}\mathrm d\beta
&=\int_0^\infty\left[\int_{\beta^\tau}^\infty\alpha^{q-q_0-1}\,\mathrm d\alpha\right]^{p_0/q_0}
\beta^{p_0-1}\lambda_f(\beta)\,\mathrm d\beta\\
&=(q_0-q)^{-p_0/q_0}\int_0^\infty\beta^{p-1}\lambda_f(\beta)\,\mathrm d\beta\\
&=|q-q_0|^{-p_0/q_0}p^{-1}\|f\|_p^p.
\end{align*}
类似的计算给出
\begin{align*}
\int_0^\infty\left[\int_0^\infty\phi_1^{q_1/p_1}\,\mathrm d\alpha\right]^{p_1/q_1}\mathrm d\beta
=|q-q_1|^{-p_1/q_1}p^{-1}\|f\|_p^p.
\end{align*}
结合\cref{align:ch06-marcinkiewicz-two-integrals,align:ch06-marcinkiewicz-minkowski}，得到
\begin{align*}
\sup_{\|f\|_p=1}\|Tf\|_q\leqslant
B_p=2q^{1/q}\left[\sum_{j=0}^1 C_j^{q_j}
\left(\frac{p_j}p\right)^{q_j/p_j}|q-q_j|^{-1}\right]^{1/q}.
\end{align*}
由 $|T(cf)|=c|Tf|$（$c>0$），可恢复一般 $f$ 的估计。原书将 $B_p$ 所述性质的验证留作容易的练习。

还须处理 $q_0=\infty$ 或 $q_1=\infty$ 的例外情形，分为三种。

\textbf{情形 I：$p_1=q_1=\infty$。}此时 $p_0\leqslant q_0<\infty$。把截断高度改为 $A=\alpha/C_1$，便有 $\|Th_A\|_\infty\leqslant C_1\|h_A\|_\infty\leqslant\alpha$，所以 $\lambda_{Th_A}(\alpha)=0$。这对应\cref{align:ch06-marcinkiewicz-two-integrals}中 $\phi_1=0$，且 $\phi_0$ 的区域边界 $\alpha^\sigma$ 换成 $\alpha/C_1$。与上面相同的论证给出
\begin{align*}
\|Tf\|_q\leqslant2\left[qC_0^{q_0}C_1^{q-q_0}
\left(\frac{p_0}p\right)^{q_0/p_0}|q-q_0|^{-1}\right]^{1/q}\|f\|_p.
\end{align*}
\textbf{情形 II：$p_0<p_1<\infty$、$q_0<q_1=\infty$。}仍选 $A$ 使 $\lambda_{Th_A}(\alpha)=0$，正确的选择为
\begin{align*}
A=(\alpha/d)^\sigma,\qquad
d=C_1\left(\frac{p_1\|f\|_p^p}p\right)^{1/p_1},\qquad
\sigma=\frac{p_1}{p_1-p}.
\end{align*}
这里 $\sigma$ 为\cref{align:ch06-marcinkiewicz-sigma}中 $q_1\to\infty$ 时的极限。因为 $p_1>p$，
\begin{align*}
\|Th_A\|_\infty^{p_1}
&\leqslant C_1^{p_1}\|h_A\|_{p_1}^{p_1}
=C_1^{p_1}p_1\int_0^A\beta^{p_1-1}\lambda_f(\beta)\,\mathrm d\beta\\
&\leqslant C_1^{p_1}p_1 A^{p_1-p}\int_0^A\beta^{p-1}\lambda_f(\beta)\,\mathrm d\beta\\
&\leqslant C_1^{p_1}\frac{p_1}p(\alpha/d)^{p_1}\|f\|_p^p=\alpha^{p_1}.
\end{align*}
与情形 I 一样，$\phi_1=0$；当 $\|f\|_p=1$ 时，含 $\phi_0$ 的积分受只依赖所述参数的常数控制，得到结论。

\textbf{情形 III：$p_0<p_1<\infty$、$q_1<q_0=\infty$。}论证与情形 II 基本相同，只须取 $A=(\alpha/d)^\sigma$，并选择 $d$ 使 $\lambda_{Tg_A}(\alpha)=0$。
\end{proof}
% SUPPLEMENT: id=v14-ch06-marcinkiewicz-completion; parent=ch06-marcinkiewicz; kind=omitted-case+proof-detail+motivation+hypotheses; origin=independent-derivation
\begin{remark}
\textbf{遗漏计算与三个无穷端点的完整补证（重构新增）。}下面把原证明的类似计算和省略分支逐项展开。可以使用严格正的弱型界 $C_j$，即使某个端点的最优常数为零，也选一个正的上界即可。

首先，截断的端点可积性不是预先假定的。因为 $p_0<p<p_1$，
\begin{align*}
\|g_A\|_{p_0}^{p_0}&\leqslant\frac{p_0}p A^{p_0-p}\|f\|_p^p,\\
\|h_A\|_{p_1}^{p_1}&\leqslant\frac{p_1}p A^{p_1-p}\|f\|_p^p
\quad(p_1<\infty),\qquad \|h_A\|_\infty\leqslant A.
\end{align*}
第一式在 $\beta\geqslant A$ 上用 $\beta^{p_0-p}\leqslant A^{p_0-p}$，第二式在 $\beta\leqslant A$ 上用 $\beta^{p_1-p}\leqslant A^{p_1-p}$，再用层积分公式。因此端点弱型估计确能应用于这些截断。

在两个输出端点都有限的情形，令 $\Delta_p=p_1^{-1}-p_0^{-1}\ne0$、$\Delta_q=q_1^{-1}-q_0^{-1}\ne0$。原定义给出 $p^{-1}-p_0^{-1}=t\Delta_p$、$q^{-1}-q_0^{-1}=t\Delta_q$，在另一端同为因子 $t-1$，这证明\cref{align:ch06-marcinkiewicz-sigma}的中间等式。其共同值非零，且
\begin{align*}
\tau\frac{p_j(q-q_j)}{q_j}=p-p_j.
\end{align*}
对尚未展开的 $j=1$：若 $q_1>q_0$，则 $\sigma>0$、$q-q_1<0$，区域 $\beta<\alpha^\sigma$ 等价于 $\alpha>\beta^\tau$，故内积分为
\begin{align*}
\int_{\beta^\tau}^\infty\alpha^{q-q_1-1}\,\mathrm d\alpha
=\frac{\beta^{\tau(q-q_1)}}{q_1-q}.
\end{align*}
若 $q_1<q_0$，则 $\sigma<0$、$q-q_1>0$，同一原区域等价于 $\alpha<\beta^\tau$，内积分为
\begin{align*}
\int_0^{\beta^\tau}\alpha^{q-q_1-1}\,\mathrm d\alpha
=\frac{\beta^{\tau(q-q_1)}}{q-q_1}.
\end{align*}
两种情况取 $p_1/q_1$ 次幂并乘 $\beta^{p_1-1}\lambda_f(\beta)$ 后，指数均为 $p-1$；再积分即得原证明所用的 $|q-q_1|^{-p_1/q_1}p^{-1}\|f\|_p^p$。所有函数非负，允许积分先为无穷；Minkowski 后所得有限值才证明输出属于 $L^q$，没有预先使用结论。

\textbf{求和系数的原文核对。}原式 (6.33) 前两项从弱型界直接产生 $C_0^{q_0},C_1^{q_1}$；把两项合写为求和时，既未换元也未重新定义常数，故仍须为 $C_j^{q_j}$。原页此处的 $C_j^{p_j}$ 与前行一般不相等，而原书下一页 $B_p$ 的公式又恢复为 $C_j^{q_j}$。正文保留了这个原文差异的明确说明，并使用代数上一致的系数。

\textbf{情形 I 的积分。}先归一化 $\|f\|_p=1$。这里 $q/p=q_0/p_0$，且在 $g_A$ 的区域 $\beta>\alpha/C_1$ 上有 $0<\alpha<C_1\beta$。Minkowski 后的内积分给出
\begin{align*}
\left[\int_0^{C_1\beta}\alpha^{q-q_0-1}\,\mathrm d\alpha\right]^{p_0/q_0}
=\frac{C_1^{(q-q_0)p_0/q_0}\beta^{(q-q_0)p_0/q_0}}
{(q-q_0)^{p_0/q_0}}.
\end{align*}
乘上 $\beta^{p_0-1}\lambda_f(\beta)$，指数变成 $p-1$；外积分为 $1/p$，所以得到正文所列常数。

\textbf{情形 II 的积分。}仍先设 $\|f\|_p=1$，记 $d=C_1(p_1/p)^{1/p_1}$、$\sigma=p_1/(p_1-p)>0$。已有 $\|Th_A\|_\infty\leqslant\alpha$。剩余区域为
\begin{align*}
\beta>(\alpha/d)^\sigma\quad\Longleftrightarrow\quad
0<\alpha<d\beta^{1/\sigma}.
\end{align*}
于是内积分等于 $d^{q-q_0}\beta^{(q-q_0)/\sigma}/(q-q_0)$。插值关系给出 $p_0+(q-q_0)p_0/(\sigma q_0)=p$，故重复外积分后可取
\begin{align*}
B_p=2\left[\frac{qC_0^{q_0}(p_0/p)^{q_0/p_0}d^{q-q_0}}{q-q_0}\right]^{1/q}.
\end{align*}
一般非零 $f$ 由正齐次性缩放回来；零函数由端点界处理，毋须在 $d=0$ 时除法。

\textbf{情形 III 的全部计算。}先设 $\|f\|_p=1$，取
\begin{align*}
\sigma=\frac{p_0}{p_0-p}<0,\qquad d=C_0(p_0/p)^{1/p_0},\qquad A=(\alpha/d)^\sigma.
\end{align*}
由尾部截断界，
\begin{align*}
\|Tg_A\|_\infty^{p_0}
\leqslant C_0^{p_0}\frac{p_0}p A^{p_0-p}
=C_0^{p_0}\frac{p_0}p(\alpha/d)^{p_0}=\alpha^{p_0}.
\end{align*}
因此只有 $Th_A$ 项留下。因 $\sigma<0$，其区域满足
\begin{align*}
\beta<(\alpha/d)^\sigma\quad\Longleftrightarrow\quad
0<\alpha<d\beta^{1/\sigma}.
\end{align*}
此时 $q>q_1$，故内积分为 $d^{q-q_1}\beta^{(q-q_1)/\sigma}/(q-q_1)$。取 $p_1/q_1$ 次幂、乘 $\beta^{p_1-1}\lambda_f(\beta)$，再用 $p_1+(q-q_1)p_1/(\sigma q_1)=p$，得到外积分 $d^{(q-q_1)p_1/q_1}(q-q_1)^{-p_1/q_1}/p$。所以
\begin{align*}
B_p=2\left[\frac{qC_1^{q_1}(p_1/p)^{q_1/p_1}d^{q-q_1}}{q-q_1}\right]^{1/q}
\end{align*}
可用。正齐次性再次恢复一般输入，完成三个分支。

\textbf{常数在端点的性质。}先设 $p_0<p_1$。若两个 $q_j$ 都有限，由插值的仿射关系，$q$ 是 $p$ 的光滑函数，且在两个有限端点有非零有限导数：从 $1/q=a+b/p$ 得 $\mathrm dq/\mathrm dp=bq^2/p^2$，这里 $b=\Delta_q/\Delta_p\ne0$。故 $|q-q_j|$ 与 $|p-p_j|$ 在端点附近互相控制。其余正因子有有限界，因而 $B_p\leqslant C(1+|p-p_j|^{-1/q})\leqslant C'(1+|p-p_j|^{-1})$，最后一步在 $|p-p_j|<1$ 时用 $q\geqslant1$。

情形 I 在有限端点同理；当 $p\to\infty$ 时 $q=(q_0/p_0)p\to\infty$。常数的对数中，$\log q/q$、$\log(q-q_0)/q$、$\log p/q$ 均趋零，而 $(q-q_0)\log C_1/q\to\log C_1$，所以 $B_p\to2C_1$。情形 II 在 $p\to p_0$ 时仍为有限端点估计；当 $p\to p_1$ 时 $q\to\infty$、$d\to C_1$，同一对数计算给 $B_p\to2C_1$。情形 III 在 $p\to p_1$ 时为有限端点估计，在 $p\to p_0$ 时 $q\to\infty$、$d\to C_0$，故 $B_p\to2C_0$。这些界蕴含原书声称的全部端点增长控制。

\textbf{相同输入指数的记号。}若 $p_0=p_1$，则 $p$ 固定，而 $q$ 仍随 $t$ 改变；这时原定理的 $B_p$ 应理解为固定插值点的常数，更明确可记为 $B_{p,q}$。不能声称常数只依赖这个固定的 $p$ 而对所有中间 $q$ 一致；习题中的分层积分给出它对 $q$ 的实际依赖。原书关于 $p\to p_j$ 的增长断言适用于 $p_0\ne p_1$ 的非退化插值区间。

\textbf{动机与逐条件说明。}截断高度随输出高度改变，让大值部分使用低输入指数、小值部分使用高输入指数，再让两个积分都汇成 $\int\beta^{p-1}\lambda_f$。$p_j\leqslant q_j$ 恰保证交换积分时的 Minkowski 指数 $q_j/p_j\geqslant1$；$q_0\ne q_1$ 提供两个不同尾部幂，使中间 $q$ 两侧的积分各自收敛。$0<t<1$ 保证严格处于这些收敛区间，并使 $q<\infty$；$p_j\geqslant1$ 还保证 $(\beta-A)^{p_0-1}\leqslant\beta^{p_0-1}$。次线性只用来分割水平集与缩放，不要求复线性；输入和空间使截断可交给同一个 $T$。测度空间本身无需 $\sigma$ 有限：交换的两次积分都在正半轴的 Lebesgue 测度上，原空间上的层积分公式已对任意测度建立。
\end{remark}
% END SUPPLEMENT: id=v14-ch06-marcinkiewicz-completion
% SOURCE p=221; printed=206; object=ch06-marcinkiewicz-special-cases-request; region=top
这个证明中的长公式看起来可能令人却步，但思想相当简单。为使其更清楚，原书建议读者把两个特殊而重要的情形写出来：(i) $p_0=q_0=1$、$p_1=q_1=2$；(ii) $p_0=q_0=1$、$p_1=q_1=\infty$。
% SUPPLEMENT: id=v14-ch06-marcinkiewicz-special-cases; parent=ch06-marcinkiewicz-special-cases-request; kind=omitted-proof+motivation+hypotheses; origin=independent-derivation
\begin{remark}
\textbf{两个指定特例的完整计算（重构新增）。}第一种情形取 $1<p<2$，此时输入、输出的中间指数相同。对每个 $\alpha>0$ 用 $A=\alpha$ 截断。弱端点界及层积分给出
\begin{align*}
\|Tf\|_p^p
\leqslant2^pp\left[
C_0\int_0^\infty\alpha^{p-2}\int_\alpha^\infty\lambda_f(\beta)\,\mathrm d\beta\,\mathrm d\alpha
+2C_1^2\int_0^\infty\alpha^{p-3}\int_0^\alpha\beta\lambda_f(\beta)\,\mathrm d\beta\,\mathrm d\alpha
\right].
\end{align*}
由非负积分的 Tonelli 定理，第一个二重积分为
\begin{align*}
\int_0^\infty\lambda_f(\beta)\int_0^\beta\alpha^{p-2}\,\mathrm d\alpha\,\mathrm d\beta
=\frac{\|f\|_p^p}{p(p-1)};
\end{align*}
第二个二重积分为
\begin{align*}
\int_0^\infty\beta\lambda_f(\beta)\int_\beta^\infty\alpha^{p-3}\,\mathrm d\alpha\,\mathrm d\beta
=\frac{\|f\|_p^p}{p(2-p)}.
\end{align*}
因此 $\|Tf\|_p\leqslant2[C_0/(p-1)+2C_1^2/(2-p)]^{1/p}\|f\|_p$。$p>1$ 保证第一内积分在零端收敛，$p<2$ 保证第二内积分在无穷端收敛。

第二种情形取 $1<p<\infty$，令 $A=\alpha/C_1$。因为 $\|Th_A\|_\infty\leqslant\alpha$，
\begin{align*}
\|Tf\|_p^p
&\leqslant2^ppC_0\int_0^\infty\alpha^{p-2}
\int_{\alpha/C_1}^\infty\lambda_f(\beta)\,\mathrm d\beta\,\mathrm d\alpha\\
&=2^ppC_0\int_0^\infty\lambda_f(\beta)
\int_0^{C_1\beta}\alpha^{p-2}\,\mathrm d\alpha\,\mathrm d\beta\\
&=\frac{2^pC_0C_1^{p-1}}{p-1}\|f\|_p^p.
\end{align*}
这里层积分的因子 $p$ 在最后一步抵消，故 $\|Tf\|_p\leqslant2[C_0C_1^{p-1}/(p-1)]^{1/p}\|f\|_p$。这个特例解释了下一推论中 $p\downarrow1$ 时出现的增长；无穷输出端点直接消掉一项，故不再有有限的上端点限制。
\end{remark}
% END SUPPLEMENT: id=v14-ch06-marcinkiewicz-special-cases
% SOURCE p=221; printed=206; object=ch06-interpolation-comparison; region=upper-middle
比较这两个插值定理：Marcinkiewicz 定理要求某些 Riesz–Thorin 定理不要求的指数限制；不过原书所关心的应用满足这些限制。除此之外，它的假设更弱：允许 $T$ 次线性而非线性，也只要求两个端点的弱型估计。两者都给出 $L^p(\mu)$ 到 $L^q(\nu)$ 的有界性，但 Riesz–Thorin 定理给出的算子范数界精确得多。因此这两个定理互不包含。
\end{document}
